%% =====================================================================
%%  第3章 関連研究
%%  書くこと：関連する研究の整理と，本研究との違い（本研究の位置づけ）
%%  章の最後で既存研究の課題を示し，次の「提案手法」につなげる．
%% =====================================================================
\chapter{関連研究}
\label{chap:related-work}

本章では，本研究に関連する研究を整理し，本研究の位置づけを明らかにする．
\begin{writingnote}
文献を 1 つずつ紹介するのではなく，いくつかの観点（問題の設定，解法の考え方など）で分類し，観点ごとに節を分けて，それぞれの考え方・長所・限界を述べるとよい．
個々の文献については，何を対象に，どのような方法で，何を明らかにしたのかを書き，本研究との共通点と相違点を述べる．
既存研究の限界は，その研究の目的や前提を踏まえて公平に，根拠を示して指摘する．
引用する文献は自分で読んで内容を確かめ，ほかの文献での紹介だけを頼りに引用すること（孫引き）は避ける．
\end{writingnote}

\section{○○に関する研究}
\label{sec:related-a}

%% 文献の引用：\cite{キー}．キーは references.bib の各文献の 1 行目に書いた名前．
%% 前の章で定義した用語や記号は，ここでそのまま使える（例：\cref{sec:problem} の問題）．
\cref{sec:problem}で述べた問題に対しては，……などの解法が研究されてきた．
例えば，Dantzig~\cite{dantzig1957discrete} は……．
また，Kellerer ら~\cite{kellerer2004knapsack} は……．

\section{△△に関する研究}
\label{sec:related-b}

%% 日本語の文献も同じように引用できる．
西~\cite{nishi2020} は……．また，西ら~\cite{nishi2011} は……．
%% 複数の文献をまとめて引用すると，番号は自動で並べ替えられ，連番は [1–3] のようにまとめられる．
このように，……に関する研究が進められている~\cite{nishi2011,nishi2020,dantzig1957discrete}．

\section{本研究の位置づけ}
\label{sec:position}

\begin{writingnote}
前の節までの整理を踏まえて，既存研究で解決されていない課題を明確に述べ，本研究がそれにどう取り組むのかを示して「提案手法」の章につなげる．
ここで示す課題は，「はじめに」の章で述べた課題と一致させる．
比較の観点は，本研究に有利なものを選ぶのではなく，問題を解くうえで重要なものを選ぶ．
「〜を扱った研究はない」と書くときは，十分に調べたうえで「筆者の知る限り」のように範囲を限って書く．
\end{writingnote}

%% 既存研究との違いは，観点を決めて表にまとめると分かりやすい．
関連研究と本研究の比較を\cref{tab:comparison}に示す．
既存研究では……が考慮されていないのに対し，本研究では……を考慮する点が異なる．

\begin{table}[tbp]
  \centering
  \caption[関連研究と本研究の比較（書式例）]{関連研究と本研究の比較（書式例．○・×は例示であり，引用文献の実際の評価ではない）}
  \label{tab:comparison}
  \begin{tabular}{lccc}
    \toprule
    & 観点 A & 観点 B & 観点 C \\
    \midrule
    文献~\cite{dantzig1957discrete} & ○ & × & × \\
    文献~\cite{nishi2011}           & ○ & ○ & × \\
    本研究                           & ○ & ○ & ○ \\
    \bottomrule
  \end{tabular}
\end{table}
